\documentclass[10pt,a4paper]{amsart}
\usepackage{amsmath,amssymb,amsthm,graphicx}
% For commutative diagrams we use Xy-pic
\usepackage[all]{xy}

% We use 2cell for 2-commutative diagrams.
\xyoption{2cell}
% Theorem environments.
%
\theoremstyle{plain}
\newtheorem{theorem}[subsection]{Theorem}
\newtheorem{proposition}[subsection]{Proposition}
\newtheorem{lemma}[subsection]{Lemma}

\theoremstyle{definition}
\newtheorem{definition}[subsection]{Definition}
\newtheorem{example}[subsection]{Example}
\newtheorem{exercise}[subsection]{Exercise}
\newtheorem{situation}[subsection]{Situation}

\theoremstyle{remark}
\newtheorem{remark}[subsection]{Remark}
\newtheorem{remarks}[subsection]{Remarks}

\numberwithin{equation}{subsection}

% Macros
%
\def\lim{\mathop{\mathrm{lim}}\nolimits}
\def\colim{\mathop{\mathrm{colim}}\nolimits}
\def\Spec{\mathop{\mathrm{Spec}}}
\def\Hom{\mathop{\mathrm{Hom}}\nolimits}
\def\Ext{\mathop{\mathrm{Ext}}\nolimits}
\def\SheafHom{\mathop{\mathcal{H}\!\mathit{om}}\nolimits}
\def\SheafExt{\mathop{\mathcal{E}\!\mathit{xt}}\nolimits}
\def\Sch{\mathit{Sch}}
\def\Mor{\mathop{\mathrm{Mor}}\nolimits}
\def\Ob{\mathop{\mathrm{Ob}}\nolimits}
\def\Sh{\mathop{\mathit{Sh}}\nolimits}
\def\NL{\mathop{N\!L}\nolimits}
\def\CH{\mathop{\mathrm{CH}}\nolimits}
\def\proetale{{pro\text{-}\acute{e}tale}}
\def\etale{{\acute{e}tale}}
\def\QCoh{\mathit{QCoh}}
\def\Ker{\mathop{\mathrm{Ker}}}
\def\Im{\mathop{\mathrm{Im}}}
\def\Coker{\mathop{\mathrm{Coker}}}
\def\Coim{\mathop{\mathrm{Coim}}}

% Boxtimes
%
\DeclareMathSymbol{\boxtimes}{\mathbin}{AMSa}{"02}

%
% Macros for moduli stacks/spaces
%
\def\QCohstack{\mathcal{QC}\!\mathit{oh}}
\def\Cohstack{\mathcal{C}\!\mathit{oh}}
\def\Spacesstack{\mathcal{S}\!\mathit{paces}}
\def\Quotfunctor{\mathrm{Quot}}
\def\Hilbfunctor{\mathrm{Hilb}}
\def\Curvesstack{\mathcal{C}\!\mathit{urves}}
\def\Polarizedstack{\mathcal{P}\!\mathit{olarized}}
\def\Complexesstack{\mathcal{C}\!\mathit{omplexes}}
% \Pic is the operator that assigns to X its picard group, usage \Pic(X)
% \Picardstack_{X/B} denotes the Picard stack of X over B
% \Picardfunctor_{X/B} denotes the Picard functor of X over B
\def\Pic{\mathop{\mathrm{Pic}}\nolimits}
\def\Picardstack{\mathcal{P}\!\mathit{ic}}
\def\Picardfunctor{\mathrm{Pic}}
\def\Deformationcategory{\mathcal{D}\!\mathit{ef}}
\usepackage{cleveref}
\crefname{originalproposition}{original proposition}{original propositions}
\Crefname{originalproposition}{Original proposition}{Original propositions}
\makeatletter
\newlabel{etale-proposition-properties-sections}{{024T}{source}{}{sourcepointer-1}{}}
\newlabel{etale-proposition-properties-sections@cref}{{[originalproposition][0][]024T}{[1][0][]source}{}{}{}}
\makeatother
\title{Realising a morphism between the closed fibres by a correspondence in the big etale site}
\author{The Stacks Project Authors}
\date{stacks-project:etale-cohomology.tex:local-snapshot-20261011; GNU FDL 1.2 or later}
\begin{document}
\maketitle
\noindent\textit{Editorial note.} All delivered mathematical text is the author's own, copied from the pinned Stacks Project chapter file etale-cohomology.tex; no statement or proof step was rewritten, reordered or removed.\par\medskip
\noindent\textit{Editorial note.} The proof is the author's complete argument: the openness of the graph and the choice of W are both inside the counted span.\par\medskip
\noindent\textit{Editorial note.} The label lemma-closed-immersion-almost-full occurs in no other staged Stacks chapter, so no cross-chapter restatement had to be weighed for this card.\par\medskip
\noindent\textit{Editorial note.} The single labelled reference occurring inside the counted spans points to a chapter of the Project not staged in this card's source tree; it is carried in assembly.external\_references as a non-load-bearing navigation pointer recorded at the invoking line.\par\medskip
\noindent\textit{Editorial note.} The preamble is the chapter's own native preamble excerpt already used for this source: the Xy-pic setup (native preamble lines 14-18) and the native theorem declarations with the mathematical macros (native preamble lines 154-217).\par\medskip
\noindent\textit{Allowed standard prerequisites (editorial).}\begin{itemize}
\item Scheme U\_Z = U x\_X Z attached to a scheme U etale over X and a closed immersion i : Z -> X. Same chapter, the notation U\_Z for the base change of an object of X\_etale along the closed immersion. Supplies the source and target of the given morphism h and the closed subscheme M\_Z of the product.
\item Graph of a morphism and its openness for an etale projection. Same chapter together with the chapter on Etale Morphisms, the fact that the graph of a morphism to an etale scheme is the image of a section, hence open. Supplies the openness of the graph Gamma\_h inside M\_Z, which is the crux of the proof.
\item Open subscheme with prescribed closed fibre. Same chapter, the standard fact that an open subscheme of the closed fibre M\_Z can be lifted to an open subscheme W of M. Supplies the open subscheme W of M = U x\_X U' whose closed fibre is Gamma\_h.
\end{itemize}\medskip
\ifdefined\hypertarget\hypertarget{sourcepointer-1}{}\fi
\noindent\textit{Source reference.} Original proposition 024T: sources/stacks-etale-cohomology/pinned-original-source.tex, source lines 6614--6614. Cross-chapter navigation pointer invoked at proof lines 6612-6614, where the author uses the properties of sections of an etale morphism, in particular that the image of a section is open. The target chapter is not staged in this card's source tree, so the pointer records the invoking line; it only resolves the reference for the reader, and no bibliography citation is involved.\par\medskip
\begin{lemma}
\label{lemma-closed-immersion-almost-full}
Let $i : Z \to X$ be a closed immersion of schemes.
Let $U, U'$ be schemes \'etale over $X$. Let $h : U_Z \to U'_Z$
be a morphism over $Z$. Then there exists a diagram
$$
\xymatrix{
U & W \ar[l]_a \ar[r]^b & U'
}
$$
in $X_\etale$ such that $a_Z : W_Z \to U_Z$
is an isomorphism and $h = b_Z \circ (a_Z)^{-1}$.
\end{lemma}

\begin{proof}
Consider the scheme $M = U \times_X U'$. The graph $\Gamma_h \subset M_Z$
of $h$ is open. This is true for example as $\Gamma_h$ is the image of a
section of the \'etale morphism $\text{pr}_{1, Z} : M_Z \to U_Z$, see
\'Etale Morphisms, Proposition \ref{etale-proposition-properties-sections}.
Hence there exists an open subscheme $W \subset M$ whose intersection with
the closed subset $M_Z$ is $\Gamma_h$. Set $a = \text{pr}_1|_W$
and $b = \text{pr}_2|_W$.
\end{proof}

\end{document}
